\begin{solapartat}
\punts{0,2} Per calcular la massa de Júpiter ens cal primer la velocitat lineal o angular: $v = \dfrac{2\pi r}{T}$ o $\omega = \dfrac{2\pi}{T}$

\punts{0,1} $v = \dfrac{2\pi r}{T} = 1,73\times 10^{4}\ \mathrm{m/s}$ o $\omega = \dfrac{2\pi}{T} = 4,11\times 10^{-5}\ \mathrm{rad/s}$.

\punts{0,1} La segona llei de Newton estableix que: $\vec{F} = M_{Io}\vec{a}$

Segons la llei de gravitació universal, el mòdul de la força s'expressa com:

$F = G\,\dfrac{M_{Io}M_J}{r^2}$ \quad \punts{0,2}

Per tant obtenim que: $a = G\,\dfrac{M_J}{r^2}$ \quad \punts{0,1}

\punts{0,2} D'altra banda, considerant que Io descriu un moviment circular uniforme al voltant de Júpiter, la seva acceleració centrípeta és: $a = \dfrac{v^2}{r}$ o $a = \omega^2 r$,

I la massa de Júpiter és

\punts{0,25} $M_J = \dfrac{r\cdot v^2}{G} = 1,90\times 10^{27}\ \mathrm{kg}$ o $M_J = \dfrac{r^3\cdot\omega^2}{G} = 1,90\times 10^{27}\ \mathrm{kg}$ \quad \punts{0,1}
\end{solapartat}

\begin{solapartat}
\punts{1} De la tercera llei de Kepler:
\[ \frac{T_{Io}^2}{r_{Io}^3} = \frac{T_{Cal\cdot listo}^2}{r_{Cal\cdot listo}^3} \Longrightarrow r_{Cal\cdot listo} = r_{Io}\left(\frac{T_{Cal\cdot listo}}{T_{Io}}\right)^{\frac{2}{3}} = 1,88\times 10^{9}\ \mathrm{m} \quad \text{\punts{0,25}} \]

\destaca{Alternativament}

\punts{0,1}, La velocitat angular és: o $\omega = \dfrac{2\pi}{T} = 4,36\times 10^{-6}\ \mathrm{rad/s}$ .

\punts{0,25} $G\,\dfrac{M_{Cal\cdot listo}M_J}{r^2} = M_{Cal\cdot listo}\,\omega^2 r$

\punts{0,8} $r = \sqrt[3]{G\,\dfrac{M_J}{\omega^2}} = 1,88\times 10^{9}\ \mathrm{m}$ \quad \punts{0,1}
\end{solapartat}
